% =============================================================================
%   Sistema Embarcado de Navegação Semi-Assistida para Cadeiras de Rodas
%   Ponto de Controle 4 (PC4) — Sistemas Operacionais Embarcados — UnB
%
%   Compilação:  pdflatex pc4.tex && pdflatex pc4.tex
%
%   Auto-contido: bibliografia inline (\bibitem) e diagramas em TikZ.
%   Para usar figuras externas (fotos, screenshots da GUI/RViz), substitua os
%   blocos TikZ pelos \includegraphics correspondentes.
% =============================================================================

\documentclass[conference]{IEEEtran}
\IEEEoverridecommandlockouts

\usepackage[utf8]{inputenc}
\usepackage[T1]{fontenc}
\usepackage[brazilian]{babel}
\usepackage{graphicx}
\usepackage{amsmath,amssymb}
\usepackage{cite}
\usepackage{url}
\usepackage{xcolor}
\usepackage{tikz}
\usepackage{listings}
\usetikzlibrary{shapes.geometric, arrows.meta, positioning, fit, backgrounds, calc}

% Estilo discreto para trechos curtos de código
\lstset{
  basicstyle=\ttfamily\footnotesize,
  breaklines=true,
  columns=fullflexible,
  keepspaces=true,
  frame=single,
  framesep=3pt,
  xleftmargin=2pt,
  xrightmargin=2pt
}

\begin{document}

\title{Sistema Embarcado de Navegação Semi-Assistida\\
para Cadeiras de Rodas Motorizadas:\\
Supervisor de Segurança Semi-Assistido Baseado em LiDAR}

\author{
\IEEEauthorblockN{Vítor Gonçalves de Andrade Silva — 261025182}
\IEEEauthorblockA{Faculdade de Ciências e Tecnologia em Engenharia\\
Engenharia Eletrônica\\
St. Leste Projeção A Gama Leste, Brasília -- DF, 72444-240\\
Email: 261025182@aluno.unb.br}
\and
\IEEEauthorblockN{Thiago Henrique Oliveira Souza — 180131672}
\IEEEauthorblockA{Faculdade de Ciências e Tecnologia em Engenharia\\
Engenharia Eletrônica\\
St. Leste Projeção A Gama Leste, Brasília -- DF, 72444-240\\
Email: 180131672@aluno.unb.br}
}

\maketitle

% =============================================================================
\begin{abstract}
Este Ponto de Controle~4 (PC4) consolida a missão central do projeto: a
implementação efetiva do supervisor de segurança semi-assistido baseado em
LiDAR, planejado arquiteturalmente no PC3. Foi estendido o protocolo serial
entre a Raspberry Pi~5 e o microcontrolador, acrescentando um comando de
condução assistida (\texttt{drive\_cmd}) e um portão de segurança
(\texttt{drive\_cfg}) com \emph{watchdog} temporizado, de modo que a perda de
qualquer enlace resulta em parada segura, nunca em movimento inesperado. Sobre
essa base, o supervisor foi realizado de duas formas complementares: (i)~uma
camada ROS\,2 (\texttt{wheelchair\_ros}) com um nó de controle compartilhado
que avalia uma \emph{função custo} sobre direções candidatas obtidas do
\texttt{/scan}, decidindo entre frear, desviar ou parar; e (ii)~uma versão
pragmática integrada à interface gráfica de bancada, que modula em tempo real
o teto de \emph{duty} dos motores conforme a folga frontal, com zonas de
redução e parada. Adicionalmente, foi diagnosticada a causa-raiz das falhas
intermitentes do LiDAR observadas no PC3 --- subtensão do barramento USB da
Raspberry Pi~5 sob pico de corrente --- e corrigido um defeito de exaustão de
memória que travava o sistema operacional após alguns minutos de operação. Os
ensaios de bancada validam a parada automática diante de obstáculos e os
mecanismos de \emph{failsafe}, cumprindo o princípio de que o sistema de alto
nível pode \emph{frear}, mas nunca \emph{acelerar}.
\end{abstract}

% =============================================================================
\section{Introdução}

O projeto desenvolve um sistema de direção semi-assistida para cadeiras de
rodas motorizadas, voltado a usuários com deficiência motora severa, baixo
controle de tronco ou ausência de controle motor fino~\cite{fehr2000,branco2018}.
Comandos acidentais e espasmos no joystick podem produzir acelerações
repentinas e colisões; a literatura clínica aponta que parcela expressiva dos
pacientes é considerada inapta ao uso de cadeiras motorizadas justamente pela
inadequação das interfaces de controle~\cite{fehr2000}. A proposta converte a
cadeira em um veículo semi-autônomo, dividido em um modelo mestre--escravo: a
Raspberry Pi~5, sob Ubuntu~24.04 e ROS\,2 Jazzy~\cite{ros2docs}, executa a
percepção e a assistência de alto nível, enquanto o microcontrolador mantém o
laço de tempo real do joystick e dos motores~\cite{ogata2010}.

Ao final do PC3, a arquitetura distribuída sobre ROS\,2 estava integrada, o IMU
publicava a 250\,Hz e o pacote próprio servia de espinha dorsal ao sistema;
restavam, contudo, três pendências críticas: a indisponibilidade do LiDAR
durante os testes, a ausência de telemetria útil das rodas e --- sobretudo ---
a não implementação do supervisor de segurança, então apenas definido em
princípio. O PC4 ataca diretamente essa lacuna central. Suas contribuições são:

\begin{enumerate}
\item \textbf{Protocolo de controle assistido}: extensão do firmware com os
comandos \texttt{drive\_cfg} (portão de segurança: armar/desarmar, teto de
\emph{duty} e rampas) e \texttt{drive\_cmd} (condução assistida por roda),
ambos sob \emph{watchdog}, com estados explícitos de operação reportados na
telemetria.

\item \textbf{Supervisor de segurança por função custo}: nó ROS\,2 que assina
o \texttt{/scan}, avalia direções candidatas por uma função custo que pondera
obstáculo e desvio, e decide entre reduzir velocidade, desviar lateralmente ou
parar o avanço.

\item \textbf{Parada automática de bancada}: implementação direta na interface
gráfica de controle, que modula o teto de \emph{duty} em função da folga no
cone frontal, com zonas de redução e parada parametrizáveis, usada para a
demonstração física.

\item \textbf{Robustez do sistema embarcado}: diagnóstico do problema de
energia USB da Raspberry Pi~5 (subtensão sob pico de corrente) e correção de um
defeito de exaustão de memória que travava o sistema operacional, tema
diretamente associado a temporização e tolerância a falhas.
\end{enumerate}

% =============================================================================
\section{Evolução da Arquitetura}

A divisão mestre--escravo do PC3 foi preservada. A novidade do PC4 está no
\emph{caminho de comando assistido}: além de receber a telemetria do
microcontrolador, a Raspberry Pi passa a poder \emph{intervir} no movimento de
forma limitada e auditável. A Fig.~\ref{fig:fluxo} ilustra o fluxo de dados do
supervisor. O joystick físico permanece como fonte primária de intenção; o
supervisor só pode \emph{atenuar} ou \emph{bloquear} o avanço, nunca ampliá-lo.

% --- Fluxo de dados do supervisor -------------------------------------------
\begin{figure*}[!t]
\centering
\resizebox{0.92\textwidth}{!}{%
\begin{tikzpicture}[
  font=\footnotesize,
  every node/.style={align=center},
  blk/.style={draw, rounded corners=2pt, minimum height=10mm,
              minimum width=30mm, inner sep=4pt, fill=gray!6},
  sens/.style={draw, rounded corners=2pt, minimum height=10mm,
               minimum width=24mm, inner sep=4pt, fill=gray!14},
  arr/.style={-Latex, thick},
  lab/.style={font=\scriptsize, fill=white, inner sep=1pt}
]
\node[sens] (esp) at (0,0) {ESP32 / firmware\\joystick + motores};
\node[blk]  (bridge) at (4.2,0) {\texttt{esp\_bridge}\\(ponte serial)};
\node[blk]  (shared) at (9.2,0) {\texttt{shared\_control}\\(função custo)};
\node[sens] (lidar) at (9.2,-2.0) {LiDAR\\\texttt{/scan}};

\draw[arr] (esp.east) -- (bridge.west)
  node[midway,above,lab] {telemetria};
\draw[arr] (bridge.east) -- (shared.west)
  node[midway,above,lab] {\texttt{/joystick\_cmd\_vel}};
\draw[arr] (lidar.north) -- (shared.south)
  node[midway,right,lab] {\texttt{/scan}};

% Retorno: cmd_vel -> bridge -> drive_cmd -> esp
\draw[arr] (shared.north) -- ++(0,0.9) -| (bridge.north)
  node[pos=0.25,above,lab] {\texttt{/cmd\_vel} (atenuado/desviado)};
\draw[arr] (bridge.south) -- ++(0,-1.4) -| (esp.south)
  node[pos=0.25,below,lab] {\texttt{drive\_cfg} + \texttt{drive\_cmd}};
\end{tikzpicture}%
}
\caption{Fluxo de dados do supervisor semi-assistido. A intenção do joystick,
lida pelo firmware, é republicada pela ponte em \texttt{/joystick\_cmd\_vel}.
O nó \texttt{shared\_control} combina essa intenção com o \texttt{/scan} e
publica um \texttt{/cmd\_vel} possivelmente atenuado ou desviado, que a ponte
converte em \texttt{drive\_cmd} para o firmware. A iniciativa de avanço é
sempre do usuário; o supervisor apenas freia ou desvia.}
\label{fig:fluxo}
\end{figure*}

\subsection{Protocolo de controle assistido e estados de operação}

O firmware passou a aceitar três mensagens JSON sobre a serial USB:

\begin{lstlisting}[basicstyle=\ttfamily\scriptsize]
{"type":"drive_cfg","armed":true,"max_duty":0.20,
 "accel":1.5,"decel":3.0}
{"type":"drive_cmd","left":0.0,"right":0.0}
{"type":"stop"}
\end{lstlisting}

O \texttt{drive\_cfg} funciona como \emph{portão de segurança}: define se o
estágio de potência está armado e impõe o teto de \emph{duty} e as rampas de
aceleração/desaceleração. O \texttt{drive\_cmd} carrega o comando assistido
normalizado por roda, em $[-1,1]$. O firmware seleciona, a cada ciclo, entre
três modos de condução, reportados no campo \texttt{drive\_mode} da telemetria
e ilustrados na Fig.~\ref{fig:estados}:

\begin{itemize}
\item \textbf{manual}: não há \texttt{drive\_cmd} recente; o firmware aplica a
mixagem diferencial do joystick físico, saturada pelo teto de \texttt{drive\_cfg};
\item \textbf{assist}: há \texttt{drive\_cmd} recente; o firmware usa os comandos
de roda calculados pela Raspberry Pi;
\item \textbf{assist\_timeout}: havia condução assistida, mas o
\texttt{drive\_cmd} expirou (perda do nó de controle); o firmware \emph{para},
sem retornar silenciosamente ao modo manual.
\end{itemize}

Dois \emph{watchdogs} independentes garantem segurança por temporização: a
ausência de \texttt{drive\_cfg} fresco desarma o estágio de potência (queda da
ponte ou do USB), e a expiração do \texttt{drive\_cmd} no estado assistido leva
à parada. Assim, qualquer falha de comunicação degrada para um estado seguro.

% --- Máquina de estados do firmware -----------------------------------------
\begin{figure}[!t]
\centering
\resizebox{0.95\columnwidth}{!}{%
\begin{tikzpicture}[
  font=\footnotesize,
  node distance=15mm and 18mm,
  st/.style={draw, rounded corners=2pt, minimum height=9mm,
             minimum width=22mm, align=center, fill=gray!8},
  arr/.style={-Latex, thick},
  lab/.style={font=\scriptsize, fill=white, inner sep=1pt}
]
\node[st, fill=gray!18] (dis) {desarmado};
\node[st, right=of dis]  (man) {manual};
\node[st, right=of man]  (ass) {assist};
\node[st, below=of ass, fill=red!12] (to) {assist\_timeout\\(parado)};

\draw[arr] (dis.east) -- (man.west)
  node[midway,above,lab] {\texttt{drive\_cfg}\\armado};
\draw[arr] (man.east) -- (ass.west)
  node[midway,above,lab] {\texttt{drive\_cmd}};
\draw[arr] (ass.south) -- (to.north)
  node[midway,right,lab] {\texttt{drive\_cmd}\\expira};
\draw[arr] (to.west) -- ++(-1.0,0) |- (man.south)
  node[pos=0.75,below,lab] {novo \texttt{drive\_cmd} / manual};
\draw[arr] (man.north) -- ++(0,0.9) -| (dis.north)
  node[pos=0.25,above,lab] {\texttt{stop} / \texttt{drive\_cfg} expira};
\end{tikzpicture}%
}
\caption{Estados de condução do firmware. Perdas de enlace (USB, ponte ou nó de
controle) levam sempre a estados seguros: desarmado ou parado.}
\label{fig:estados}
\end{figure}

\subsection{Supervisor de segurança por função custo}

O nó \texttt{shared\_control} é o núcleo da assistência. A cada ciclo
($\sim$20\,Hz) ele recebe a intenção do usuário $(v_u,\dot\psi_u)$ e a varredura
do LiDAR, e avalia um conjunto de direções candidatas
$\delta \in [-\delta_{\max}, \delta_{\max}]$ em torno da frente. Para cada
candidata define-se a \emph{folga} $c(\delta)$ como a menor distância medida
dentro de um cone de meia-abertura $\theta_c$ centrado em $\delta$:
\begin{equation}
c(\delta) = \min_{\,|\angle(\theta_i,\delta)| \le \theta_c}\; r_i,
\end{equation}
\noindent onde $(r_i,\theta_i)$ são os pares distância/ângulo do \texttt{/scan}.
O termo de obstáculo penaliza folgas pequenas de forma contínua entre os
limiares de parada $d_{\text{stop}}$ e de redução $d_{\text{slow}}$:
\begin{equation}
g(c) =
\begin{cases}
0, & c \ge d_{\text{slow}},\\[2pt]
\dfrac{d_{\text{slow}} - c}{\,d_{\text{slow}} - d_{\text{stop}}\,}, & d_{\text{stop}} < c < d_{\text{slow}},\\[8pt]
C_{\text{blk}}, & c \le d_{\text{stop}}.
\end{cases}
\end{equation}
A função custo combina obstáculo e desvio em relação à direção desejada
$\delta_c = \mathrm{clamp}(\dot\psi_u,-\delta_{\max},\delta_{\max})$:
\begin{equation}
J(\delta) = w_o\,g\!\big(c(\delta)\big)
          + w_d\,\frac{|\delta - \delta_c|}{\delta_{\max}}.
\end{equation}
A direção escolhida é $\delta^\star = \arg\min_\delta J(\delta)$. A decisão de
arbitragem é:
\begin{itemize}
\item se \emph{todas} as candidatas estão bloqueadas ($g=C_{\text{blk}}$) ou
$c(\delta^\star) \le d_{\text{stop}}$: \textbf{parar} o avanço;
\item caso a frente esteja totalmente livre ($c(0) \ge d_{\text{slow}}$):
\textbf{nenhuma intervenção};
\item caso intermediário: reduzir a velocidade pelo fator
$s\!\big(c(\delta^\star)\big)$ e aplicar uma correção limitada de esterçamento.
\end{itemize}
O fator de redução $s(\cdot)$ é linear na faixa $[d_{\text{stop}},d_{\text{slow}}]$,
e os comandos de saída são
\begin{align}
v_{\text{out}}      &= v_u \cdot s\!\big(c(\delta^\star)\big),\\
\dot\psi_{\text{out}} &= \dot\psi_u + K_{\text{assist}}\,\delta^\star,
\end{align}
\noindent com $K_{\text{assist}}$ ajustável (em particular, $K_{\text{assist}}=0$
desativa o desvio, restando apenas freio/parada --- modo conservador usado nos
primeiros ensaios). Essa formulação por setores e custo direcional aproxima-se
da família de métodos de histograma de campo vetorial~\cite{borenstein1991},
adaptada ao requisito de \emph{preservação da intenção do usuário}.

A ponte \texttt{esp\_bridge} converte $(v_{\text{out}},\dot\psi_{\text{out}})$ em
comandos normalizados de roda por mixagem diferencial e os transmite como
\texttt{drive\_cmd}, mantendo o \texttt{drive\_cfg} vivo como portão de
segurança. Por padrão a ponte nasce \emph{desarmada}: só transmite movimento
quando explicitamente armada.

\subsection{Zonas de segurança e parada automática de bancada}

Para a demonstração física foi implementada também uma versão direta na
interface gráfica de controle, sem depender da pilha completa de controle
ROS\,2. Como o firmware já satura o joystick pelo teto de \texttt{drive\_cfg} e
a interface já transmite esse teto periodicamente, a parada automática foi
obtida \emph{modulando o teto de \texttt{max\_duty} em tempo real} conforme a
folga no cone frontal. A lógica restringe a leitura a um cone de meia-abertura
$\theta_c=30^\circ$ à frente e ignora retornos a menos de $0{,}15$\,m (estrutura
da própria cadeira), conforme as zonas da Fig.~\ref{fig:zonas}:

\begin{itemize}
\item \textbf{zona vermelha} ($c \le 0{,}60$\,m): teto de \emph{duty} forçado a
zero --- \emph{parada};
\item \textbf{zona amarela} ($0{,}60 < c \le 1{,}10$\,m): teto reduzido
linearmente com a folga;
\item \textbf{livre} ($c > 1{,}10$\,m): teto nominal.
\end{itemize}

\begin{figure}[!t]
\centering
\resizebox{0.82\columnwidth}{!}{%
\begin{tikzpicture}[font=\footnotesize]
% Zona amarela (reducao) - cone frontal +-30 graus, raio 1.10 (escala 2.5)
\fill[orange!50, opacity=0.25] (0,0) -- (60:2.75) arc[start angle=60, end angle=120, radius=2.75] -- cycle;
% Zona vermelha (parada) - mesmo cone, raio 0.60 (escala 2.5 -> 1.5)
\fill[red!60, opacity=0.45] (0,0) -- (60:1.5) arc[start angle=60, end angle=120, radius=1.5] -- cycle;
\draw[orange!70!black, thin] (0,0) -- (60:2.75);
\draw[orange!70!black, thin] (0,0) -- (120:2.75);
\draw[orange!70!black, thin] (60:2.75) arc[start angle=60, end angle=120, radius=2.75];
\draw[red!70!black, thin] (60:1.5) arc[start angle=60, end angle=120, radius=1.5];

% Cadeira e LiDAR
\filldraw[fill=gray!25, draw=gray!60] (-0.6,-0.7) rectangle (0.6,0.4);
\node[gray!40!black] at (0,-0.15) {\scriptsize cadeira};
\fill[gray!70!black] (0,0.4) circle (0.10);
\node[right=0.5mm of {(0.10,0.4)}, gray!40!black] {\scriptsize LiDAR};

% Cotas
\node[red!60!black] at (90:1.0) {\scriptsize $0{,}60$\,m};
\node[orange!60!black] at (90:2.3) {\scriptsize $1{,}10$\,m};
\node[gray!50!black] at (123:1.9) {\scriptsize $\pm 30^\circ$};

% Obstaculo de exemplo na zona vermelha
\filldraw[fill=gray!60, draw=gray!80] (0.15,1.2) rectangle (0.5,1.55);
\draw[gray!50, ->] (1.3,1.4) -- (0.55,1.4) node[midway,above,gray!40!black]
  {\scriptsize obstáculo};

% Legenda
\node[anchor=west] at (1.6,0.4) {\scriptsize zona amarela: reduz $v_{\max}$};
\node[anchor=west] at (1.6,0.0) {\scriptsize zona vermelha: parada};
\end{tikzpicture}%
}
\caption{Zonas implementadas na parada automática de bancada, restritas ao cone
frontal de $\pm 30^\circ$. A zona amarela atenua o teto de \emph{duty}
linearmente com a folga; a zona vermelha ($\le 0{,}60$\,m) impõe parada.
Retornos a menos de $0{,}15$\,m são descartados como estrutura da cadeira.}
\label{fig:zonas}
\end{figure}

\subsection{Resiliência: temporização e exaustão de recursos}

A interface de controle concentra três produtores de dados assíncronos (ESP,
LiDAR e IMU) drenados por um laço de renderização periódico. Verificou-se que o
reagendamento desse laço ocorria \emph{após} todas as rotinas de desenho;
qualquer exceção transitória em uma delas interrompia o reagendamento e, por
consequência, a drenagem da fila de eventos. Como as \emph{threads} produtoras
seguiam enfileirando (o LiDAR a 10\,Hz, com $\sim$1500 pontos por varredura), a
memória crescia de forma ilimitada até travar todo o sistema operacional após
alguns minutos. A correção aplicou dois princípios clássicos de sistemas
embarcados:

\begin{enumerate}
\item \textbf{laço auto-resiliente}: o reagendamento passou a ocorrer em
cláusula \texttt{finally}, de modo que uma falha de desenho é registrada mas
\emph{nunca} interrompe a drenagem;
\item \textbf{fila limitada com descarte}: a fila de eventos recebeu teto e os
produtores de alto volume passaram a descartar amostras antigas quando ela
satura (apenas a varredura mais recente é útil), eliminando a possibilidade de
crescimento ilimitado de memória.
\end{enumerate}

% =============================================================================
\section{Resultados Experimentais}

Os ensaios do PC4 foram conduzidos no hardware final, em bancada, com as rodas
suspensas em todos os passos que podem acionar motor. Adotou-se sempre teto
conservador ($\texttt{max\_duty} \le 0{,}20$) e início sem desvio
($K_{\text{assist}}=0$).

\subsection{Experimento 1: Energia USB da Raspberry Pi~5 e recuperação do LiDAR}

A indisponibilidade do LiDAR ao final do PC3 foi reinvestigada. O registro do
\emph{kernel} revelou desconexões e re-enumerações simultâneas de múltiplos
dispositivos USB --- inclusive do próprio LiDAR --- padrão típico de
\emph{brownout} do barramento. A Raspberry Pi~5 limita a corrente total das
portas USB quando não alimentada pela fonte de 27\,W; com o motor do LiDAR e
os demais conversores seriais drenando corrente, o barramento sofria subtensão
momentânea, reiniciando os dispositivos. A remoção do IMU do barramento
estabilizou imediatamente o conjunto, confirmando a hipótese: a ``falha'' do
PC3 foi, em boa medida, um problema de \emph{energia}, não de sensor. A
Tabela~\ref{tab:usb} resume o diagnóstico.

\begin{table}[!t]
\caption{Diagnóstico de estabilidade do barramento USB (Raspberry Pi~5)}
\label{tab:usb}
\centering
\footnotesize
\renewcommand{\arraystretch}{1.2}
\begin{tabular}{@{}p{0.40\columnwidth}p{0.52\columnwidth}@{}}
\hline
\textbf{Observação} & \textbf{Interpretação} \\
\hline
Desconexões simultâneas de LiDAR e ESP no log do \emph{kernel} & Subtensão do barramento sob pico de corrente \\
Conjunto estabiliza ao remover o IMU & Excesso de corrente agregada \\
Portas \texttt{ttyUSB*} trocam de número entre reinícios & Enumeração não determinística \\
\hline
\end{tabular}
\end{table}

Como mitigação imediata, adotou-se o endereçamento dos dispositivos por
caminhos estáveis (\emph{symlinks} \texttt{by-id}, equivalentes às regras
\emph{udev} do PC3), eliminando a dependência da ordem de enumeração, e a
operação da demonstração sem o IMU, que não é necessário ao supervisor de
parada. Recomenda-se, para a fase de campo, alimentar o LiDAR por
\emph{hub} USB com fonte própria. Recuperada a estabilidade, o LiDAR voltou a
publicar \texttt{/scan} a $\sim$10\,Hz, desbloqueando a percepção.

\subsection{Experimento 2: Telemetria e protocolo de condução assistida}

Com o microcontrolador conectado, a telemetria foi capturada diretamente da
serial USB a 115\,200\,bps, confirmando a emissão contínua de pacotes
\texttt{drive} já com os campos do protocolo estendido --- entre eles
\texttt{max\_duty}, \texttt{accel}, \texttt{decel} e os contadores de idade dos
comandos (\texttt{cfg\_age\_ms}, \texttt{assist\_age\_ms}). A presença desses
campos comprova que o firmware em execução implementa o protocolo
\texttt{drive\_cfg}/\texttt{drive\_cmd}, condição necessária para a condução
assistida. A ponte serial ROS\,2 abriu a porta, republicou a intenção do
joystick em \texttt{/joystick\_cmd\_vel} e manteve o \texttt{drive\_cfg}
periódico, partindo no estado desarmado por padrão.

\subsection{Experimento 3: Parada automática diante de obstáculo}

Validou-se a parada automática de bancada com o LiDAR em operação e o cone
frontal alinhado à frente da cadeira. A Tabela~\ref{tab:zonas} resume o
comportamento observado ao aproximar um obstáculo da frente, em concordância
com a lógica das zonas: parada total na zona vermelha, redução proporcional na
zona amarela e teto nominal em frente livre. O indicador da interface
acompanhou as transições \textsc{livre}~$\rightarrow$~\textsc{reduz}~$\rightarrow$~\textsc{parada}
em tempo real, e o teto de \emph{duty} transmitido ao firmware foi a zero na
zona vermelha, impedindo o giro das rodas mesmo com o joystick solicitando
avanço.

\begin{table}[!t]
\caption{Resposta da parada automática à folga frontal (bancada)}
\label{tab:zonas}
\centering
\footnotesize
\renewcommand{\arraystretch}{1.2}
\begin{tabular}{@{}p{0.24\columnwidth}p{0.22\columnwidth}p{0.40\columnwidth}@{}}
\hline
\textbf{Folga frontal} & \textbf{Zona} & \textbf{Teto de \emph{duty} aplicado} \\
\hline
$\le 0{,}60$\,m            & vermelha & $0$ (parada) \\
$0{,}80$\,m                & amarela  & $\approx 40\%$ do nominal \\
$> 1{,}10$\,m / só lateral & livre    & nominal \\
$< 0{,}15$\,m              & ignorada & descartada (estrutura) \\
\hline
\end{tabular}
\end{table}

\subsection{Experimento 4: Mecanismos de \emph{failsafe}}

Os caminhos de degradação segura foram exercitados isoladamente, todos com
resultado conforme o projetado:

\begin{itemize}
\item \textbf{Perda de \texttt{/scan}} (LiDAR encerrado com pedido de avanço):
o controle bloqueia o avanço por ausência de varredura recente;
\item \textbf{Perda do nó de controle}: a ponte passa a enviar \texttt{stop} e
o firmware desarma;
\item \textbf{Perda do enlace USB/ponte}: expirado o \texttt{drive\_cfg}, o
firmware desarma e zera os motores;
\item \textbf{\texttt{drive\_cmd} expirado com \texttt{drive\_cfg} ainda
fresco}: o firmware entra em \texttt{assist\_timeout} e \emph{para}, sem
retornar sozinho ao modo manual.
\end{itemize}

\subsection{Experimento 5: Correção do travamento por exaustão de memória}

Reproduziu-se a falha de travamento e confirmou-se sua causa-raiz (interrupção
do laço de drenagem por exceção transitória, com crescimento ilimitado da fila
de eventos). Após a aplicação do laço auto-resiliente e da fila limitada com
descarte, a operação prolongada deixou de degradar: na pior hipótese, a
interface perde atualizações momentâneas, mas a memória permanece limitada e o
sistema operacional não trava. O resultado ilustra, na prática, os conceitos de
temporização determinística e tolerância a falhas centrais à disciplina.

% =============================================================================
\section{Conclusões}

O PC4 cumpriu a missão central do projeto: implementar e validar o supervisor
de segurança semi-assistido baseado em LiDAR, antes apenas previsto. O firmware
ganhou um protocolo de condução assistida com portão de segurança e
\emph{watchdogs}, garantindo que toda falha de comunicação degrade para parada
ou desarme. Sobre essa base, o supervisor foi realizado tanto como nó ROS\,2 de
controle compartilhado por função custo quanto como parada automática direta na
interface de bancada, ambos respeitando o princípio de que o sistema de alto
nível pode \emph{frear}, mas nunca \emph{acelerar}. A Tabela~\ref{tab:status}
resume a situação dos subsistemas.

\begin{table}[!t]
\caption{Situação dos principais subsistemas ao final do PC4}
\label{tab:status}
\centering
\footnotesize
\renewcommand{\arraystretch}{1.2}
\begin{tabular}{@{}p{0.24\columnwidth}p{0.24\columnwidth}p{0.40\columnwidth}@{}}
\hline
\textbf{Subsistema} & \textbf{Status} & \textbf{Observação} \\
\hline
LiDAR / \texttt{/scan}        & Operacional & Recuperado; falha do PC3 atribuída a subtensão USB. \\
Telemetria do ESP             & Operacional & Protocolo \texttt{drive\_cfg}/\texttt{drive\_cmd} em execução. \\
Supervisor (função custo)     & Implementado & Nó ROS\,2 \texttt{shared\_control} com freio/desvio/parada. \\
Parada automática             & Validada     & Em bancada, por modulação do teto de \emph{duty}. \\
\emph{Failsafes}              & Validados    & \emph{Watchdogs} de \texttt{drive\_cfg} e \texttt{drive\_cmd}. \\
Energia USB (Pi 5)            & Mitigada     & Endereçamento \texttt{by-id}; \emph{hub} alimentado recomendado. \\
Robustez da interface         & Corrigida    & Laço resiliente + fila limitada (sem travamento). \\
Desvio autônomo com motor     & A validar    & Implementado; teste motorizado em campo pendente. \\
\hline
\end{tabular}
\end{table}

Como limitações, o desvio autônomo foi implementado mas ainda não exercitado em
condução motorizada de campo, e a parada baseada em teto de \emph{duty} inibe
também a ré nas proximidades do obstáculo (comportamento seguro, porém
conservador). O bloqueio direcional --- parar apenas o avanço e preservar a ré
--- requer o uso pleno do \texttt{drive\_cmd} sobrepondo a mixagem do joystick,
caminho já disponível na camada ROS\,2.

Os próximos passos para a fase final são:
\begin{enumerate}
\item realizar o ensaio motorizado de campo do supervisor com a cadeira no solo,
partindo de teto conservador;
\item alimentar o LiDAR por \emph{hub} USB com fonte própria, eliminando o risco
de subtensão sob pico de corrente;
\item calibrar os limiares $d_{\text{stop}}$, $d_{\text{slow}}$, a abertura do
cone e o ganho $K_{\text{assist}}$ em função da velocidade máxima da cadeira e
da latência total;
\item implementar o bloqueio direcional via \texttt{drive\_cmd}, preservando a
manobra de recuo;
\item integrar a pose fundida (\texttt{/odometry/filtered}) ao supervisor para
referência de movimento mais robusta.
\end{enumerate}

% =============================================================================
\begin{thebibliography}{99}

\bibitem{fehr2000}
L.~Fehr, W.~E. Langbein e S.~B. Skaar, ``Adequacy of power wheelchair
control interfaces for persons with severe disabilities: A clinical
survey,'' \emph{Journal of Rehabilitation Research and Development},
vol.~37, no.~3, pp.~353--360, 2000.

\bibitem{borenstein1991}
J.~Borenstein e Y.~Koren, ``The Vector Field Histogram --- Fast Obstacle
Avoidance for Mobile Robots,'' \emph{IEEE Transactions on Robotics and
Automation}, vol.~7, no.~3, pp.~278--288, 1991.

\bibitem{hossain2022}
M.~A. Hossain \emph{et~al.}, ``Collision Avoidance for Autonomous
Wheelchair Using Integrated SLAM and YOLO,'' em \emph{6th International
Conference on I-SMAC (IoT in Social, Mobile, Analytics and Cloud)}, 2022.

\bibitem{ahmed2021}
S.~S.~S.~J. Ahmed \emph{et~al.}, ``SLAM Based Indoor Autonomous
Navigation System For Electric Wheelchair,'' em \emph{2021 International
Conference on ICICT4SD}, 2021.

\bibitem{ogata2010}
K.~Ogata, \emph{Engenharia de Controle Moderno}, 5\textsuperscript{a}~ed.
São Paulo, SP: Pearson Prentice Hall, 2010.

\bibitem{branco2018}
M.~D.~C.~C. Branco, ``Desenvolvimento de um protótipo de cadeira de rodas
motorizada controlada por movimentos oculares,'' Trabalho de Conclusão
de Curso, UFRN, Natal, 2018.

\bibitem{ros2docs}
Open Robotics, ``Robot Operating System 2 (ROS\,2) Jazzy Jalisco
Documentation,'' 2024. Disponível em:
\url{https://docs.ros.org/en/jazzy/}.

\bibitem{robotlocalization}
T.~Moore e D.~Stouch, ``A Generalized Extended Kalman Filter
Implementation for the Robot Operating System,'' em \emph{Proceedings of
the 13th International Conference on Intelligent Autonomous Systems},
Springer, 2014.

\bibitem{slamtec}
Slamtec, ``RPLIDAR Low Cost 360 Degree Laser Range Scanner Datasheet,''
2020.

\bibitem{witmotion}
WitMotion, ``IWT603 Industrial Grade IMU Product Specifications,'' 2024.

\end{thebibliography}

\end{document}
